\section{Introduction}
% P1: service-level problem and why interpretation belongs in the latency budget.
\IEEEPARstart{N}{atural-language} interfaces let users describe an edge service
and its execution requirements together: read the text in an image, keep the
image at its originating site, and return the result before a deadline.
Turning that request into an executable job requires both semantic
interpretation and a placement decision. In edge computing, communication,
resource availability, and placement jointly determine whether the service
responds in time~\cite{shi2016edge,mao2017survey,santos2021lowlatency}.
An interpreter on the admission path consumes part of that same response
budget, delaying execution and leaving less time to complete the service.

% P2: bounded output motivates substitution; completion is the acceptance criterion.
For a bounded service catalog, interpretation often requires only a small
set of decisions, such as service type, placement permission, quality tier,
and urgency. A large language model (LLM) can express these decisions through
a short structured response, but the downstream scheduler needs the selected
values rather than a generated explanation. A decision-oriented model could
reduce this overhead even relative to a generative model configured for
concise output. At the service level, the goal is to reduce this decision
overhead while preserving correct, on-time completion under the request's
execution requirements.

% P3: established work and the concrete comparison examined here.
Intent-based networking already separates desired outcomes from their
implementation~\cite{clemm2022intentbased}. Language-driven network systems
and service prototypes connect interpretation to configuration and
orchestration~\cite{jacobs2021hey,wang2024netconfeval,angi2025llnet,miyaoka2025chatdriven,nisiotis2026prompt}.
Small models and discriminative extraction offer further ways to specialize
bounded language tasks~\cite{belcak2025small,zaratiana2025gliner2}.
We examine Jev, a decision-oriented service accessed through an application
programming interface (API)~\cite{friedman2026jev}, as the interpreter in an
edge-service admission path. Queueing, communication, and execution all
contribute to response time; their interaction with interpretation determines
how much a faster decision benefits the service.

% P4: three answerable questions, aligned with the evaluation and Results.
We investigate three linked questions. How does Jev compare with a concise
structured-output DeepSeek deployment in decision latency, semantic
correctness, and billed API cost? Does the substitution lower full response
time while retaining correct service completion? Under what workload
conditions is faster interpretation most useful, and how does it interact
with caching for repeated requests?

% P5: methodological response and the distinct role of each evidence layer.
We address these questions through live API measurements combined with
modeled execution, followed by a real two-node optical character recognition
(OCR) service. The real service transfers images, executes recognition, and
checks the returned text against reference annotations. A fixed rule parser
and self-hosted Qwen model extend the comparison beyond the two hosted APIs.

% P6: concrete system contribution, measured service benefit, and deployment guidance.
The contributions are as follows.
\begin{enumerate}
\item We integrate Jev into an edge-service admission system through a
four-field intent contract and a shared validator and scheduler. The system
accounts for interpretation, queueing, communication, and execution on one
timeline, connecting the choice of interpreter to service deadlines and
completion.
\item We measure a latency and API-cost benefit from this substitution.
Jev reduces median decision latency by 15.9--26.5\% across three blocks.
In the real OCR service, it matches DeepSeek's correct-completion count in
seven paired conditions and exceeds it in one. Without caching, median
full response time on requests completed correctly by both systems is
11.1--25.3\% lower; API fees per correct completion are 69.0--70.6\% lower
across the eight conditions.
\item We identify how request reuse changes the value of faster decisions.
Applying the same cache policy to all backends shows that latency savings
are concentrated in fresh interpretation, while repeated-request caching
largely removes the difference. This comparison guides the use of fast
decision models alongside caching in service orchestration.
\end{enumerate}
